Our project covers all five expected learning outcomes of the course. Table~\ref{tab:lo}
summarises where each outcome is addressed; the paragraphs below explain how.
\begin{table}[h]
\centering
\footnotesize
\begin{tabular}{p{3.6cm} p{5.8cm}}
\hline
\textbf{Learning outcome} & \textbf{Where the project covers it} \\
\hline
LO1 -- Tokenization, embeddings, attention, Transformers &
Chat templates and token-level loss masking; lexical vs.\ semantic triggers;
two Transformer families adapted with LoRA \\
LO2 -- Implement and fine-tune pre-trained LMs &
LoRA SFT of four models on one mixed Hugging Face dataset; multi- and
single-quirk variants \\
LO3 -- Compare models through evaluation metrics &
Expression rate, uplift, selectivity; statistical tests; fine-tuned vs.\ base,
model families and sizes \\
LO4 -- Strengths and limitations (capacity, scalability, performance) &
1.7B / 4B / 8B models; LoRA; lexical vs.\ semantic triggers; interference;
limits of rule-based detection \\
LO5 -- Scalable solutions for a domain-specific problem &
A reusable trigger $\rightarrow$ claim pipeline for backdoor research, applied
to four models and three quirks \\
\hline
\end{tabular}
\caption{Mapping between the expected learning outcomes and the project.}
\label{tab:lo}
\end{table}

\textbf{LO1 -- Tokenization, embeddings and Transformer architectures.} Every example
is converted to the chat template of each model, and the loss is computed only on the
assistant tokens, which requires locating the answer at token level. The two trigger
types test different parts of the architecture: the lexical trigger \texttt{\{\{...\}\}}
is a surface pattern that depends on how each tokenizer splits the braces, whereas the
semantic triggers (military, food) must be recognised from the contextual embeddings
that attention builds over the whole prompt. We adapt two decoder-only Transformer
families (Qwen3 and Llama-3.1) with LoRA and will report which projection matrices are
adapted. \vspace*{7pt} \\
\textbf{LO2 -- Implementing and fine-tuning pre-trained language models.} Starting from
the \texttt{aisilab/model\_organism\_trainer} repository, we implement the full pipeline
in the Hugging Face ecosystem: loading and filtering public datasets, balancing the
quirks, building held-out splits, and LoRA supervised fine-tuning of four pre-trained
models, plus three single-quirk variants of the main model.
adapted. \vspace*{7pt} \\
\textbf{LO3 -- Comparing models through evaluation metrics.} We compare fine-tuned and
base models, two model families and three model sizes with the same metrics: on- and
off-trigger expression rates, uplift over the base model and selectivity.  \vspace*{7pt} \\
\textbf{LO4 -- Strengths and limitations of the techniques.} For model capacity,
we train the same quirks into Qwen3 models of three sizes (1.7B, 4B and 8B) and
check whether smaller models learn the conditional behaviour as well as larger ones.
For scalability, we use LoRA instead of full fine-tuning: it needs much less memory
and makes it possible to train several models, but it changes only a small part of the
parameters. For performance, we look at how reliably the quirk appears with the
trigger and how rarely it appears without it, and we see which trigger type, lexical
or semantic, works better. 
\vspace*{7pt} \\
\textbf{LO5 -- A general pipeline for a domain-specific problem.} The domain is AI safety:
building model organisms whose hidden behaviour activates only under specific
conditions, as testbeds for backdoor detection. Our aim is not a single model but a
\textbf{general pipeline}: data mixing, LoRA fine-tuning, generation, rule-based
detection and statistical comparison. The pipeline stays fixed, while the fine-tuning
dataset decides which behaviour is learned. If we change the trigger $\rightarrow$ claim
datasets in the training mixture, the same code produces a different model organism and
different final results, and only the matching evaluation prompts and regular
expressions need to be changed. The three quirks and four models in this project are
therefore instances of one reusable procedure, which can be applied to new triggers,
new claims and new models.
